% TikZ 图：玻尔氢原子模型与能级
% Day 27：近代物理与原子物理入门
\begin{tikzpicture}[scale=0.88,>=Latex,line cap=round,line join=round,
  every node/.style={font=\small}]
  % Bohr orbit picture
  \begin{scope}[shift={(-3.15,0)}]
    \fill[red!75!black] (0,0) circle (2.4pt);
    \node[below=4pt] at (0,-0.05) {原子核};
    \foreach \r in {0.55,1.05,1.55}{
      \draw[blue!65!black,thick] (0,0) circle (\r);
    }
    \node[blue!65!black,anchor=west] at (0.86,-0.26) {$n=1$};
    \node[blue!65!black,anchor=west] at (1.34,0.62) {$n=2$};
    \node[blue!65!black,anchor=west] at (1.74,1.20) {$n=3$};
    \fill[black] (0,0.55) circle (1.8pt);
    \draw[->,orange!85!black,very thick] (0,0.64) -- (0,1.48);
    \node[orange!85!black,anchor=east,align=right,font=\scriptsize,
      fill=white,inner sep=1pt] at (-0.12,1.05)
      {吸收光子\\$n=1\rightarrow n=3$};
  \end{scope}

  % Energy level picture
  \begin{scope}[shift={(1.85,-0.05)}]
    \draw[->] (-0.45,-2.15) -- (-0.45,2.20) node[above] {$E$};
    \foreach \y/\n/\e in {-1.75/1/-13.6,-0.72/2/-3.40,-0.12/3/-1.51,0.40/4/-0.85,1.25/\infty/0}{
      \draw[thick] (0,\y) -- (3.05,\y);
      \node[left=8pt] at (0,\y) {$n=\n$};
      \node[right=12pt,font=\scriptsize] at (3.05,\y) {$\e\ \mathrm{eV}$};
    }
    \draw[->,red!80!black,thick] (0.82,-0.12) -- (0.82,-0.72)
      node[midway,right,fill=white,inner sep=1pt] {放出光子};
    \draw[->,blue!75!black,thick] (2.00,-1.75) -- (2.00,1.25)
      node[midway,right,align=left,font=\scriptsize,fill=white,inner sep=1pt]
      {吸收 $13.6\ \mathrm{eV}$\\发生电离};
    \node[font=\scriptsize,gray] at (1.55,-2.17) {能级间距为示意，未按真实比例绘制};
  \end{scope}
\end{tikzpicture}
